\section{Justificación e importancia}

Esta sección explica por qué conviene medir la coherencia OpenAPI--BIAN \textbf{antes} de integrar.
No repite las definiciones del §1.1 ni la casuística del §\ref{sec:problematica}.

En arquitectura bancaria, cada \textit{release} puede desplazar la interfaz REST respecto al modelo
de negocio acordado. Farzi (2021) y Bhatia (2022) muestran que, sin trazabilidad entre capacidades
BIAN y contratos publicados, la deriva se acumula entre equipos de producto, integración y
gobernanza.

Medir esa deriva solo con revisiones manuales no escala cuando coexisten decenas de APIs y socios
externos. El plan propone cuantificar la coherencia mediante \AlignmentScore{} y las preguntas
P1--P4 (§\ref{sec:planteamiento}), de modo que una revisión documental pueda priorizarse con
evidencia objetiva (Shabbir et al., 2025).

La literatura confirma la madurez de OpenAPI como artefacto técnico (Casas et al., 2021), pero la
intersección sistemática OpenAPI--BIAN sigue delgada. Este estudio aporta el procedimiento
reproducible y la matriz de consistencia (Anexo~C) que cierra esa brecha metodológica.

El diseño descansa en representación del conocimiento, alineación de esquemas y similitud semántica
---las mismas líneas que operacionalizan S y el \textit{matching} estructural E---. Deliberadamente
no sustituye testing \textit{runtime} ni certificación regulatoria (eje E5, §\ref{sec:estado-arte}):
evalúa coherencia \textbf{documental} antes del despliegue.
